% The following code checks for the maximum Fitting length of 2-nilpotent subgroups of $GL(4,3)$. It was used for Lemma 3.1 and Theorem 3.2.

G := GL(4,3);
subs := List(ConjugacyClassesSubgroups(G), Representative);
twoNilpotentSubs := Filtered(subs, H -> IsPNilpotent(H, 2));
maxFL := Maximum(List(twoNilpotentSubs, FittingLength));
Print(maxFL);


% The following code checks for the Fitting length of 2-nilpotent exceptional cases of $(n,q) = (4,7)$ in \cite{Foulser}. It was used for Theorem 3.2.

LoadPackage("primgrp");

deg := 7^4;
ord := 7^4 * 2^7 * 3 * 5;

cands := AllPrimitiveGroups(
    NrMovedPoints, deg,
    Size, ord,
    IsSolvableGroup, true
);

good := Filtered(cands, function(G)
    local lens;
    lens := List(Orbits(Stabilizer(G, 1), [1..deg]), Length);
    Sort(lens);
    return lens = [1, 480, 1920];
end);

for G in good do
    G0 := Stabilizer(G, 1);

    Print(
        "Size(G0) = ", Size(G0),
        ", 2-nilpotent = ", IsPNilpotent(G0, 2),
        ", Fitting length = ", FittingLength(G0),
        "\n"
    );
od;


% The following code checks for the Fitting length of 3-nilpotent exceptional cases of $(n,q) = (4,7)$ in \cite{Foulser}. It was used for Theorem 4.2.

LoadPackage("primgrp");

deg := 7^4;
ord := 7^4 * 2^7 * 3 * 5;

cands := AllPrimitiveGroups(
    NrMovedPoints, deg,
    Size, ord,
    IsSolvableGroup, true
);

Print("Number of candidates: ", Length(cands), "\n");

good := Filtered(cands, function(G)
    local lens;
    lens := List(Orbits(Stabilizer(G, 1), [1..deg]), Length);
    Sort(lens);
    return lens = [1, 480, 1920];
end);

Print("Number with subdegrees 1,480,1920: ", Length(good), "\n");

for G in good do
    G0 := Stabilizer(G, 1);

    Print(
        "Size(G0) = ", Size(G0),
        ", 3-nilpotent = ", IsPNilpotent(G0, 3),
        ", Fitting length = ", FittingLength(G0),
        "\n"
    );
od;

% The following code was used for Theorem 3.4 and Theorem 4.4 to generate the tables.

OnWordByPositions := function(w, g)
    local u, i;
    u := [];

    for i in [1..Length(w)] do
        u[i^g] := w[i];
    od;

    return u;
end;

BaseWordOfType := function(pi)
    local w, color, k, i;

    w := [];
    color := 1;

    for k in pi do
        for i in [1..k] do
            Add(w, color);
        od;
        color := color + 1;
    od;

    return w;
end;

Fpi := function(pi)
    local vals, v, prod;

    vals := Set(pi);
    prod := 1;

    for v in vals do
        prod := prod * Factorial(Length(Filtered(pi, x -> x = v)));
    od;

    return prod;
end;

PiString := function(pi)
    local str, i;

    str := "$(";

    for i in [1..Length(pi)] do
        str := Concatenation(str, String(pi[i]));

        if i < Length(pi) then
            str := Concatenation(str, ",");
        fi;
    od;

    str := Concatenation(str, ")$");

    return str;
end;

Ntriv := function(G, pi)
    local base, words, seen, count, w, orb;

    base := BaseWordOfType(pi);
    words := Set(PermutationsList(base));

    seen := [];
    count := 0;

    for w in words do
        if not w in seen then
            orb := Set(List(Elements(G), g -> OnWordByPositions(w, g)));

            if Size(Stabilizer(G, w, OnWordByPositions)) = 1 then
                count := count + 1;
            fi;

            seen := Union(seen, orb);
        fi;
    od;

    return count;
end;

PrintNtrivLatexRows := function(G, m)
    local parts, pi, n;

    parts := List(Partitions(m), p -> Reversed(p));

    for pi in parts do
        n := Ntriv(G, pi);

        if n <> 0 then
            Print(Length(pi), " & ",
                  PiString(pi), " & ",
                  n, " & ",
                  Fpi(pi), " \\\\\n");
        fi;
    od;
end;

S3 := SymmetricGroup(3);
PrintNtrivLatexRows(S3, 3);

F10 := Group(
    (1,2,3,4,5),
    (2,5)(3,4)
);
PrintNtrivLatexRows(F10, 5);

F20 := Group(
    (1,2,3,4,5),
    (2,3,5,4)
);
PrintNtrivLatexRows(F20, 5);

F42 := Group(
    (1,2,3,4,5,6,7),
    (2,4,3,7,5,6)
);
PrintNtrivLatexRows(F42, 7);

t1 := (1,2,3)(4,5,6)(7,8,9);
t2 := (1,4,7)(2,5,8)(3,6,9);
rD8 := (2,4,3,7)(5,6,9,8);
sD8 := (4,7)(5,8)(6,9);
C3xC3_D8 := Group(
    t1, t2,
    rD8, sD8
);
PrintNtrivLatexRows(C3xC3_D8, 9);

t1 := (1,2,3)(4,5,6)(7,8,9);
t2 := (1,4,7)(2,5,8)(3,6,9);
aSD := (2,4,5,8,3,7,9,6);
bSD := (2,5)(3,9)(4,7);
C3xC3_SD16 := Group(
    t1, t2,
    aSD, bSD
);
PrintNtrivLatexRows(C3xC3_SD16, 9);

A4 := AlternatingGroup(4);
PrintNtrivLatexRows(A4, 4);

trans8 := (1,2)(3,5)(4,8)(6,7);
mult8  := (2,3,4,5,6,7,8);
frob8  := (3,4,6)(5,8,7);
AGamma23 := Group(
    trans8,
    mult8,
    frob8
);
PrintNtrivLatexRows(AGamma23, 8);